% Options for packages loaded elsewhere
\PassOptionsToPackage{unicode}{hyperref}
\PassOptionsToPackage{hyphens}{url}
%
\documentclass[
  ignorenonframetext,
  aspectratio=169,
]{beamer}
\usepackage{pgfpages}
\setbeamertemplate{caption}[numbered]
\setbeamertemplate{caption label separator}{: }
\setbeamercolor{caption name}{fg=normal text.fg}
\beamertemplatenavigationsymbolshorizontal
% Prevent slide breaks in the middle of a paragraph
\widowpenalties 1 10000
\raggedbottom
\setbeamertemplate{part page}{
  \centering
  \begin{beamercolorbox}[sep=16pt,center]{part title}
    \usebeamerfont{part title}\insertpart\par
  \end{beamercolorbox}
}
\setbeamertemplate{section page}{
  \centering
  \begin{beamercolorbox}[sep=12pt,center]{part title}
    \usebeamerfont{section title}\insertsection\par
  \end{beamercolorbox}
}
\setbeamertemplate{subsection page}{
  \centering
  \begin{beamercolorbox}[sep=8pt,center]{part title}
    \usebeamerfont{subsection title}\insertsubsection\par
  \end{beamercolorbox}
}
\AtBeginPart{
  \frame{\partpage}
}
\AtBeginSection{
  \ifbibliography
  \else
    \frame{\sectionpage}
  \fi
}
\AtBeginSubsection{
  \frame{\subsectionpage}
}

\usepackage{amsmath,amssymb}
\usepackage{iftex}
\ifPDFTeX
  \usepackage[T1]{fontenc}
  \usepackage[utf8]{inputenc}
  \usepackage{textcomp} % provide euro and other symbols
\else % if luatex or xetex
  \usepackage{unicode-math}
  \defaultfontfeatures{Scale=MatchLowercase}
  \defaultfontfeatures[\rmfamily]{Ligatures=TeX,Scale=1}
\fi
\usepackage{lmodern}
\ifPDFTeX\else  
    % xetex/luatex font selection
\fi
% Use upquote if available, for straight quotes in verbatim environments
\IfFileExists{upquote.sty}{\usepackage{upquote}}{}
\IfFileExists{microtype.sty}{% use microtype if available
  \usepackage[]{microtype}
  \UseMicrotypeSet[protrusion]{basicmath} % disable protrusion for tt fonts
}{}
\usepackage{xcolor}
\newif\ifbibliography
\setlength{\emergencystretch}{3em} % prevent overfull lines
\setcounter{secnumdepth}{-\maxdimen} % remove section numbering

\usepackage{color}
\usepackage{fancyvrb}
\newcommand{\VerbBar}{|}
\newcommand{\VERB}{\Verb[commandchars=\\\{\}]}
\DefineVerbatimEnvironment{Highlighting}{Verbatim}{commandchars=\\\{\}}
% Add ',fontsize=\small' for more characters per line
\usepackage{framed}
\definecolor{shadecolor}{RGB}{241,243,245}
\newenvironment{Shaded}{\begin{snugshade}}{\end{snugshade}}
\newcommand{\AlertTok}[1]{\textcolor[rgb]{0.68,0.00,0.00}{#1}}
\newcommand{\AnnotationTok}[1]{\textcolor[rgb]{0.37,0.37,0.37}{#1}}
\newcommand{\AttributeTok}[1]{\textcolor[rgb]{0.40,0.45,0.13}{#1}}
\newcommand{\BaseNTok}[1]{\textcolor[rgb]{0.68,0.00,0.00}{#1}}
\newcommand{\BuiltInTok}[1]{\textcolor[rgb]{0.00,0.23,0.31}{#1}}
\newcommand{\CharTok}[1]{\textcolor[rgb]{0.13,0.47,0.30}{#1}}
\newcommand{\CommentTok}[1]{\textcolor[rgb]{0.37,0.37,0.37}{#1}}
\newcommand{\CommentVarTok}[1]{\textcolor[rgb]{0.37,0.37,0.37}{\textit{#1}}}
\newcommand{\ConstantTok}[1]{\textcolor[rgb]{0.56,0.35,0.01}{#1}}
\newcommand{\ControlFlowTok}[1]{\textcolor[rgb]{0.00,0.23,0.31}{\textbf{#1}}}
\newcommand{\DataTypeTok}[1]{\textcolor[rgb]{0.68,0.00,0.00}{#1}}
\newcommand{\DecValTok}[1]{\textcolor[rgb]{0.68,0.00,0.00}{#1}}
\newcommand{\DocumentationTok}[1]{\textcolor[rgb]{0.37,0.37,0.37}{\textit{#1}}}
\newcommand{\ErrorTok}[1]{\textcolor[rgb]{0.68,0.00,0.00}{#1}}
\newcommand{\ExtensionTok}[1]{\textcolor[rgb]{0.00,0.23,0.31}{#1}}
\newcommand{\FloatTok}[1]{\textcolor[rgb]{0.68,0.00,0.00}{#1}}
\newcommand{\FunctionTok}[1]{\textcolor[rgb]{0.28,0.35,0.67}{#1}}
\newcommand{\ImportTok}[1]{\textcolor[rgb]{0.00,0.46,0.62}{#1}}
\newcommand{\InformationTok}[1]{\textcolor[rgb]{0.37,0.37,0.37}{#1}}
\newcommand{\KeywordTok}[1]{\textcolor[rgb]{0.00,0.23,0.31}{\textbf{#1}}}
\newcommand{\NormalTok}[1]{\textcolor[rgb]{0.00,0.23,0.31}{#1}}
\newcommand{\OperatorTok}[1]{\textcolor[rgb]{0.37,0.37,0.37}{#1}}
\newcommand{\OtherTok}[1]{\textcolor[rgb]{0.00,0.23,0.31}{#1}}
\newcommand{\PreprocessorTok}[1]{\textcolor[rgb]{0.68,0.00,0.00}{#1}}
\newcommand{\RegionMarkerTok}[1]{\textcolor[rgb]{0.00,0.23,0.31}{#1}}
\newcommand{\SpecialCharTok}[1]{\textcolor[rgb]{0.37,0.37,0.37}{#1}}
\newcommand{\SpecialStringTok}[1]{\textcolor[rgb]{0.13,0.47,0.30}{#1}}
\newcommand{\StringTok}[1]{\textcolor[rgb]{0.13,0.47,0.30}{#1}}
\newcommand{\VariableTok}[1]{\textcolor[rgb]{0.07,0.07,0.07}{#1}}
\newcommand{\VerbatimStringTok}[1]{\textcolor[rgb]{0.13,0.47,0.30}{#1}}
\newcommand{\WarningTok}[1]{\textcolor[rgb]{0.37,0.37,0.37}{\textit{#1}}}

\providecommand{\tightlist}{%
  \setlength{\itemsep}{0pt}\setlength{\parskip}{0pt}}\usepackage{longtable,booktabs,array}
\usepackage{calc} % for calculating minipage widths
\usepackage{caption}
% Make caption package work with longtable
\makeatletter
\def\fnum@table{\tablename~\thetable}
\makeatother
\usepackage{graphicx}
\makeatletter
\def\maxwidth{\ifdim\Gin@nat@width>\linewidth\linewidth\else\Gin@nat@width\fi}
\def\maxheight{\ifdim\Gin@nat@height>\textheight\textheight\else\Gin@nat@height\fi}
\makeatother
% Scale images if necessary, so that they will not overflow the page
% margins by default, and it is still possible to overwrite the defaults
% using explicit options in \includegraphics[width, height, ...]{}
\setkeys{Gin}{width=\maxwidth,height=\maxheight,keepaspectratio}
% Set default figure placement to htbp
\makeatletter
\def\fps@figure{htbp}
\makeatother

\IfFileExists{plex-otf.sty}{
  %% Full TeXlive
  % \usepackage[%
  %   % math,
  %   RM={Scale=0.94},SS={Scale=0.94},SScon={Scale=0.94},TT={Scale=MatchLowercase,FakeStretch=0.9},DefaultFeatures={Ligatures=Common}
  % ]{plex-otf}
}{
  %% TinyTeX
  \usepackage{libertine}
}

%%% Load theme
% https://deic.uab.cat/~iblanes/beamer_gallery/
\IfFileExists{beamerthemegotham.sty}{
  %% Full TeXlive
  \usetheme{gotham}
  \gothamset{
    numbering=totalpagenumber,
    parttocframe default=off,
    sectiontocframe default=off,
    subsectiontocframe default=off,
  }
}{
  %% TinyTeX
  \usetheme{Madrid}
}
\makeatletter
\@ifpackageloaded{caption}{}{\usepackage{caption}}
\AtBeginDocument{%
\ifdefined\contentsname
  \renewcommand*\contentsname{Содержание}
\else
  \newcommand\contentsname{Содержание}
\fi
\ifdefined\listfigurename
  \renewcommand*\listfigurename{Список иллюстраций}
\else
  \newcommand\listfigurename{Список иллюстраций}
\fi
\ifdefined\listtablename
  \renewcommand*\listtablename{Список таблиц}
\else
  \newcommand\listtablename{Список таблиц}
\fi
\ifdefined\figurename
  \renewcommand*\figurename{Рисунок}
\else
  \newcommand\figurename{Рисунок}
\fi
\ifdefined\tablename
  \renewcommand*\tablename{Таблица}
\else
  \newcommand\tablename{Таблица}
\fi
}
\@ifpackageloaded{float}{}{\usepackage{float}}
\floatstyle{ruled}
\@ifundefined{c@chapter}{\newfloat{codelisting}{h}{lop}}{\newfloat{codelisting}{h}{lop}[chapter]}
\floatname{codelisting}{Список}
\newcommand*\listoflistings{\listof{codelisting}{Листинги}}
\makeatother
\makeatletter
\makeatother
\makeatletter
\@ifpackageloaded{caption}{}{\usepackage{caption}}
\@ifpackageloaded{subcaption}{}{\usepackage{subcaption}}
\makeatother

\ifLuaTeX
\usepackage[bidi=basic]{babel}
\else
\usepackage[bidi=default]{babel}
\fi
\babelprovide[main,import]{russian}
\babelprovide[import]{english}
% get rid of language-specific shorthands (see #6817):
\let\LanguageShortHands\languageshorthands
\def\languageshorthands#1{}
\ifLuaTeX
  \usepackage{selnolig}  % disable illegal ligatures
\fi
\usepackage{csquotes}
\usepackage{bookmark}

\IfFileExists{xurl.sty}{\usepackage{xurl}}{} % add URL line breaks if available
\urlstyle{same} % disable monospaced font for URLs
\hypersetup{
  pdftitle={Презентация 1},
  pdfauthor={Ванюшкина Татьна Валерьевна},
  pdflang={ru-RU},
  hidelinks,
  pdfcreator={LaTeX via pandoc}}


\title{Презентация 1}
\subtitle{Установка и конфигурация системы на виртуальную машину}
\author{Ванюшкина Татьна Валерьевна}
\date{2026-05-10}

\begin{document}
\frame{\titlepage}

\renewcommand*\contentsname{Содержание}
\begin{frame}[allowframebreaks]
  \frametitle{Содержание}
  \tableofcontents[hideallsubsections]
\end{frame}

\section{Информация}\label{ux438ux43dux444ux43eux440ux43cux430ux446ux438ux44f}

\begin{frame}{Докладчик}
\phantomsection\label{ux434ux43eux43aux43bux430ux434ux447ux438ux43a}
\begin{columns}[c]
\begin{column}{0.7\textwidth}
\begin{itemize}[<+->]
\tightlist
\item
  Ванюшкина Татьяна Валерьевна
\item
  д.ф.-м.н., студент
\item
  компьютерные и информационные науки
\item
  Российский университет дружбы народов им. П. Лумумбы
\item
  \href{mailto:kulyabov-ds@rudn.ru}{\nolinkurl{kulyabov-ds@rudn.ru}}
\item
  \url{https://yamadharma.github.io/ru/}
\end{itemize}
\end{column}
\end{columns}
\end{frame}

\begin{frame}{Цели и задачи}
\phantomsection\label{ux446ux435ux43bux438-ux438-ux437ux430ux434ux430ux447ux438}
\begin{itemize}[<+->]
\tightlist
\item
  Установить операционную систему на виртуальную машину
\end{itemize}
\end{frame}

\begin{frame}[fragile]{Материалы и методы}
\phantomsection\label{ux43cux430ux442ux435ux440ux438ux430ux43bux44b-ux438-ux43cux435ux442ux43eux434ux44b}
\begin{itemize}[<+->]
\tightlist
\item
  Процессор \texttt{pandoc} для входного формата Markdown
\item
  Результирующие форматы

  \begin{itemize}[<+->]
  \tightlist
  \item
    \texttt{pdf}
  \item
    \texttt{html}
  \end{itemize}
\item
  Автоматизация процесса создания: \texttt{Makefile}
\end{itemize}
\end{frame}

\section{Создание
презентации}\label{ux441ux43eux437ux434ux430ux43dux438ux435-ux43fux440ux435ux437ux435ux43dux442ux430ux446ux438ux438}

\begin{frame}{1}
\phantomsection\label{section}
\begin{itemize}[<+->]
\tightlist
\item
  загружаю в дисплейном классе операционную систему
\end{itemize}
\end{frame}

\begin{frame}{2}
\phantomsection\label{section-1}
\begin{itemize}[<+->]
\tightlist
\item
  создаю каталог с именем пользователя
\end{itemize}
\end{frame}

\begin{frame}{3}
\phantomsection\label{section-2}
\begin{itemize}[<+->]
\tightlist
\item
  запускаю виртуальную машину
\end{itemize}
\end{frame}

\begin{frame}{4}
\phantomsection\label{section-3}
\begin{itemize}[<+->]
\tightlist
\item
  создаю новую сиртуальную машину
\end{itemize}
\end{frame}

\begin{frame}{5}
\phantomsection\label{section-4}
\begin{itemize}[<+->]
\tightlist
\item
  задаю необходимые параметры
\end{itemize}
\end{frame}

\begin{frame}{6}
\phantomsection\label{section-5}
\begin{itemize}[<+->]
\tightlist
\item
  создаю пользователя
\end{itemize}
\end{frame}

\begin{frame}{7}
\phantomsection\label{section-6}
\begin{itemize}[<+->]
\tightlist
\item
  домашняя работа получаю информацию
\end{itemize}

\begin{enumerate}[<+->]
\tightlist
\item
  Версия ядра Linux (Linux version).
\item
  Частота процессора (Detected Mhz processor).
\item
  Модель процессора (CPU0).
\item
  Объем доступной оперативной памяти (Memory available).
\item
  Тип обнаруженного гипервизора (Hypervisor detected).
\item
  Тип файловой системы корневого раздела
\end{enumerate}
\end{frame}

\begin{frame}{8}
\phantomsection\label{section-7}
\begin{figure}

\centering{

\includegraphics[width=0.7\textwidth,height=\textheight]{image/1.pdf}

}

\caption{\label{fig-001}1}

\end{figure}%
\end{frame}

\begin{frame}{Процессор \texttt{pandoc}}
\phantomsection\label{ux43fux440ux43eux446ux435ux441ux441ux43eux440-pandoc}
\begin{itemize}[<+->]
\tightlist
\item
  Pandoc: преобразователь текстовых файлов
\item
  Сайт: \url{https://pandoc.org/}
\item
  Репозиторий: \url{https://github.com/jgm/pandoc}
\end{itemize}
\end{frame}

\begin{frame}[fragile]{Формат \texttt{pdf}}
\phantomsection\label{ux444ux43eux440ux43cux430ux442-pdf}
\begin{itemize}[<+->]
\tightlist
\item
  Использование LaTeX
\item
  Пакет для презентации: \href{https://ctan.org/pkg/beamer}{beamer}
\item
  Тема оформления: \texttt{metropolis}
\end{itemize}
\end{frame}

\begin{frame}[fragile]{Код для формата \texttt{pdf}}
\phantomsection\label{ux43aux43eux434-ux434ux43bux44f-ux444ux43eux440ux43cux430ux442ux430-pdf}
\begin{Shaded}
\begin{Highlighting}[]
\FunctionTok{slide\_level}\KeywordTok{:}\AttributeTok{ }\DecValTok{2}
\FunctionTok{aspectratio}\KeywordTok{:}\AttributeTok{ }\DecValTok{169}
\FunctionTok{section{-}titles}\KeywordTok{:}\AttributeTok{ }\CharTok{true}
\FunctionTok{theme}\KeywordTok{:}\AttributeTok{ metropolis}
\end{Highlighting}
\end{Shaded}
\end{frame}

\begin{frame}[fragile]{Формат \texttt{html}}
\phantomsection\label{ux444ux43eux440ux43cux430ux442-html}
\begin{itemize}[<+->]
\tightlist
\item
  Используется фреймворк \href{https://revealjs.com/}{reveal.js}
\item
  Используется \href{https://revealjs.com/themes/}{тема} \texttt{beige}
\end{itemize}
\end{frame}

\begin{frame}[fragile]{Код для формата \texttt{html}}
\phantomsection\label{ux43aux43eux434-ux434ux43bux44f-ux444ux43eux440ux43cux430ux442ux430-html}
\begin{itemize}[<+->]
\tightlist
\item
  Тема задаётся в файле \texttt{Makefile}
\end{itemize}

\begin{Shaded}
\begin{Highlighting}[]
\NormalTok{REVEALJS\_THEME = beige}
\end{Highlighting}
\end{Shaded}
\end{frame}

\begin{frame}{Результаты}
\phantomsection\label{ux440ux435ux437ux443ux43bux44cux442ux430ux442ux44b}
\begin{itemize}[<+->]
\tightlist
\item
  операционная система установлена на виртуальную машину
\end{itemize}
\end{frame}

\begin{frame}{Связь слайдов}
\phantomsection\label{ux441ux432ux44fux437ux44c-ux441ux43bux430ux439ux434ux43eux432}
\begin{itemize}[<+->]
\tightlist
\item
  Один слайд --- одна мысль
\item
  Нельзя ссылаться на объекты, находящиеся на предыдущих слайдах
  (например, на формулы)
\item
  Каждый слайд должен иметь заголовок
\end{itemize}
\end{frame}

\begin{frame}{Количество сущностей}
\phantomsection\label{ux43aux43eux43bux438ux447ux435ux441ux442ux432ux43e-ux441ux443ux449ux43dux43eux441ux442ux435ux439}
\begin{itemize}[<+->]
\tightlist
\item
  Человек может одновременно помнить \(7 \pm 2\) элемента
\item
  При размещении информации на слайде старайтесь чтобы в сумме слайд
  содержал не более 5 элементов
\item
  Можно группировать элементы так, чтобы визуально было не более 5 групп
\end{itemize}
\end{frame}

\begin{frame}{Общие рекомендации}
\phantomsection\label{ux43eux431ux449ux438ux435-ux440ux435ux43aux43eux43cux435ux43dux434ux430ux446ux438ux438}
\begin{itemize}[<+->]
\tightlist
\item
  На слайд выносится та информация, которая без зрительной опоры
  воспринимается хуже
\item
  Слайды должны дополнять или обобщать содержание выступления или его
  частей, а не дублировать его
\item
  Информация на слайдах должна быть изложена кратко, чётко и хорошо
  структурирована
\item
  Слайд не должен быть перегружен графическими изображениями и текстом
\item
  Не злоупотребляйте анимацией и переходами
\end{itemize}
\end{frame}

\begin{frame}{Представление данных}
\phantomsection\label{ux43fux440ux435ux434ux441ux442ux430ux432ux43bux435ux43dux438ux435-ux434ux430ux43dux43dux44bux445}
\begin{itemize}[<+->]
\tightlist
\item
  Лучше представить в виде схемы
\item
  Менее оптимально представить в виде рисунка, графика, таблицы
\item
  Текст используется, если все предыдущие способы отображения информации
  не подошли
\end{itemize}
\end{frame}




\end{document}
